\documentclass[10pt,a4paper]{amsart}
\usepackage[margin=19mm]{geometry}
\usepackage{amsmath,amssymb,mathtools}
\usepackage[scr=rsfs,cal=euler]{mathalfa}
\usepackage{xfrac}
\usepackage[unicode,hidelinks,bookmarks=false]{hyperref}
\newcommand{\ie}{i.e.\ }
\newcommand{\cf}{cf.\ }
\newcommand{\resp}{respectively\ }
\newcommand{\Subsection}{\subsection}
\providecommand{\nobreakdash}{\nobreak\hbox{-}}
\setlength{\emergencystretch}{2em}
\multlinegap0pt
\usepackage{amssymb}
\newtheorem{theorem}{Theorem}
\title{Spectral Bounds for Tensors Derived from Trace Functionals and Wasserstein Distance in Tensor Spaces}
\author{Hemant Sharma, Nachiketa Mishra}
\date{Source: arXiv:2605.16930v1 (CC BY 4.0)}
\begin{document}
\maketitle

\noindent\emph{Source and licence.} Spectral Bounds for Tensors Derived from Trace Functionals and Wasserstein Distance in Tensor Spaces. Version arXiv:2605.16930v1 (CC BY 4.0), distributed by arXiv under the Creative Commons Attribution 4.0 International licence, http://creativecommons.org/licenses/by/4.0/, as recorded in the arXiv licence metadata for this exact version and shown on the abstract page https://arxiv.org/abs/2605.16930v1. Full licence notice: \texttt{licenses/arxiv-2605.16930-CC-BY-4.0.txt}.\par\smallskip

\noindent\textit{Editorial scope.} Counted unit: the article's theorem theorem (source0.tex lines 575--581). The article's complete original proof of it is delivered with it (source0.tex lines 583--639, one \texttt{proof} environment of 57 lines). No mathematics is written, repaired or re-derived: the statement and the proof are the article's own lines, in the article's own notation, and no retained line is substituted or re-flowed.  No further source-local context is needed: every cross-reference of every retained line resolves inside the two delivered spans.  Nothing was omitted from any retained span, so the package carries no omission record.  No retained line contains a citation, so the wrapper carries no bibliography and no bibliography entry.  No retained line contains a figure environment, an image-inclusion command or drawing code, so no image is delivered and the package declares no asset dependency.  Editorial formatting only, in addition: an AMS \texttt{amsart} preamble that restates the article's own declarations and macros used by the retained lines, namely the environments \texttt{theorem} on separate counters, together with the standard packages those lines need.  The retained environments print in this document as Theorem 1 in order of appearance, whereas the article prints the same items with the numbering of its own full document.  The complete native source of the article as distributed in the arXiv e-print for this exact version is bundled unchanged as \texttt{sources/200755/source0.tex}.
\begin{theorem}
    For $\mathcal{A} \in \mathbb{R}^{n \times n \times p}$ fixed, if $\alpha$ and $\beta$ are any numbers satisfying
    \[
    \alpha \operatorname{tr}(\mathcal{B}) \leq \operatorname{tr}(\mathcal{A} * \mathcal{B}) \leq \beta \operatorname{tr}(\mathcal{B})
    \]
    for any positive semi-definite tensor $\mathcal{B}$, then $\alpha \leq \lambda_n(\bar{A})$ and $\lambda_1(\bar{A}) \leq \beta$, i.e., $\lambda_n(\bar{A})$ and $\lambda_1(\bar{A})$ are the tightest bounds for the inequality.
\end{theorem}

\begin{proof}
    The given inequality is
    \[
    \alpha \operatorname{tr}(\mathcal{B}) \leq \operatorname{tr}(\mathcal{A} * \mathcal{B}) \leq \beta \operatorname{tr}(\mathcal{B}),
    \]
    where $\mathcal{B}$ is any positive semi-definite (PSD) tensor. Using the $t$-product and the trace operator's properties, we write
    \[
    \operatorname{tr}(\mathcal{A} * \mathcal{B}) = \operatorname{tr}(\mathcal{A} \cdot \mathcal{B}),
    \]
    where $\mathcal{A} \cdot \mathcal{B}$ is computed in the Fourier domain as
    \[
    \widehat{\mathcal{A} * \mathcal{B}}(:, :, k) = \widehat{\mathcal{A}}(:, :, k) \cdot \widehat{\mathcal{B}}(:, :, k).
    \]

   \noindent Since $\mathcal{A}$ is Hermitian, we consider its spectral decomposition
    \[
    \widehat{\mathcal{A}}(:, :, k) = Q_k \Lambda_k Q_k^H, \quad k = 1, \ldots, p,
    \]
    where $Q_k$ is unitary, and $\Lambda_k = \operatorname{diag}(\lambda_1^{(k)}, \ldots, \lambda_n^{(k)})$ contains the eigenvalues of the $k$-th frontal slice of $\widehat{\mathcal{A}}$. The eigenvalues of $\mathcal{A}$, denoted $\lambda_1(\mathcal{A}), \ldots, \lambda_n(\mathcal{A})$, are aggregated across all slices.

  \noindent For lower bound $\alpha \leq \lambda_n(\bar{A})$,
    let $\mathcal{B} = \mathcal{U}*\mathcal{U}^H$, where $\mathcal{U}$ is a unit tensor (i.e., $\|\mathcal{U}\| = 1$), and assume $\mathcal{B}$ is aligned with the eigenvector corresponding to the smallest eigenvalue $\lambda_n(\bar{A})$ of $\mathcal{A}$. Then
    \[
    \operatorname{tr}(\mathcal{B}) = 1, \quad \operatorname{tr}(\mathcal{A} * \mathcal{B}) = \lambda_n(\bar{A}).
    \]
    Substituting into the inequality
    \[
    \alpha \leq \lambda_n(\bar{A}).
    \]

  \noindent Now, for upper bound $\lambda_1(\bar{A}) \leq \beta$,  
    similarly, let $\mathcal{B} = \mathcal{U}*\mathcal{U}^H$, where $\mathcal{U}$ is aligned with the eigenvector corresponding to the largest eigenvalue $\lambda_1(\bar{A})$ of $\mathcal{A}$. Then
    \[
    \operatorname{tr}(\mathcal{B}) = 1, \quad \operatorname{tr}(\mathcal{A} * \mathcal{B}) = \lambda_1(\bar{A}).
    \]
    Substituting into the inequality
    \[
    \lambda_1(\bar{A}) \leq \beta.
    \]
 
   \noindent To verify that $\lambda_n(\bar{A})$ and $\lambda_1(\bar{A})$ are the tightest bounds, consider any $\alpha' > \lambda_n(\bar{A})$. For $\mathcal{B}$ aligned with the smallest eigenvalue:
    \[
    \operatorname{tr}(\mathcal{A} * \mathcal{B}) = \lambda_n(\bar{A}) < \alpha' \operatorname{tr}(\mathcal{B}),
    \]
    which violates the inequality.

    \noindent Similarly, for any $\beta' < \lambda_1(\bar{A})$, there exists $\mathcal{B}$ aligned with the largest eigenvalue such that:
    \[
    \operatorname{tr}(\mathcal{A} * \mathcal{B}) = \lambda_1(\bar{A}) > \beta' \operatorname{tr}(\mathcal{B}),
    \]
    \noindent which also violates the inequality.   Hence, $\alpha = \lambda_n(\bar{A})$ and $\beta = \lambda_1(\bar{A})$ are the tightest bounds.

    \noindent The tightest bounds for the inequality are given by the smallest and largest eigenvalues of $\mathcal{A}$, satisfying
    \[
    \lambda_n(\bar{A}) \leq \frac{\operatorname{tr}(\mathcal{A} * \mathcal{B})}{\operatorname{tr}(\mathcal{B})} \leq \lambda_1(\bar{A}).
    \]
\end{proof}

\end{document}
